\subsection{\texorpdfstring{第~\ref{sec:cartan-connection}~节}{Cartan 联络}}\label{subsec:solutions-cartan-connection}

\begin{description}[font=\normalfont,leftmargin=2em,labelwidth=1.5em]
\RGsolution{ex:cartan-connection-ricci-curvature}{cartan-connection-ricci-curvature}
由曲率形式的表达式 $\Omega_\nu^\mu=\sum_{\alpha<\beta}R_{\alpha\beta\nu}{}^\mu\omega^\alpha\wedge\omega^\beta$, 可知 Riemann 张量的非零分量为
\[
  R_{j00}{}^i=-R_{0j0}{}^i=\frac{\partial^2\Phi}{\partial x^j\partial x^i}
\]
($i,j=1,2,3$). 因此 Ricci 曲率张量唯一的非零系数为
\[
  R_{00}=\sum_{i=1}^3 R_{i00}{}^i=\sum_{i=1}^3\frac{\partial^2\Phi}{\partial(x^i)^2}.
\]

\RGsolution{ex:nonflat-cartan-connection-is-not-levi-civita}{nonflat-cartan-connection-is-not-levi-civita}
Cartan 联络的 Riemann 张量的非零分量为
\[
  R_{j00}{}^i=-R_{0j0}{}^i=\frac{\partial\Phi}{\partial x^i\partial x^j}
\]
($i,j=1,2,3$). 若 Cartan 联络是 $\mathbb{R}^4$ 上某个伪 Riemann 度量的 Levi–Civita 联络, 则曲率张量的非零分量为
\[
  R_{j00\mu}=-R_{0j0\mu}
   =\sum_{i=1}^3g_{\mu i}\frac{\partial\Phi}{\partial x^i\partial x^j},
\]
其中 $(g_{\mu\nu})$ 是度量矩阵. 对伪 Riemann 度量依然成立的对称性质 $R_{j00\mu}=-R_{j0\mu0}$, 将蕴含曲率张量实际上为零 (因为 $R_{j0\mu0}=0$), 这意味着 Riemann 张量也为零. 因此, 若 Cartan 联络的曲率非零, 则它不可能是任何伪 Riemann 度量的 Levi–Civita 联络.
\end{description}
